% paper/sections/09_ccd_poisson.tex
%
% §9 圧力 Poisson 方程式の解法（§9 冒頭セクションヘッダー）
% 内容：CCD-Poisson 演算子の行列構造・変密度拡張・Balanced--Force整合性・手法比較・単体検証
%
% ═══════════════════════════════════════════════════════════════════════════════
\section{圧力 Poisson 方程式の解法：分相 PPE・HFE・欠陥補正法}
\label{sec:ppe_solve}

\noindent\textbf{本章の問い：}
二相流の圧力ステップで必要なのは，
Poisson 方程式を高次に離散化することだけではない．
投影後の速度を発散なしに戻すと同時に，
密度差による係数不連続，Young--Laplace 圧力ジャンプ，
および速度補正に使う圧力勾配を
\textbf{同じ離散面演算子の上で閉じる}必要がある．
本章はこの要請を，密度ジャンプ，圧力ジャンプ，面演算子整合，
高次片側データ，低コスト求解の五つに分解し，
それぞれに必要な数理構成を順に定める．

\begin{tcolorbox}[mybox,title={§9 の構成原理}]
\begin{enumerate}[noitemsep]
  \item \textbf{密度ジャンプ問題}：
        smoothed-Heaviside 一括 PPE では係数不連続が界面帯に残るため，
        高密度比では相内定数密度 PPE に分ける．
  \item \textbf{表面張力ジャンプ問題}：
        CSF 体積力ではなく，
        $j_{gl}=p_g-p_l=-\sigma\kappa_{lg}$ を界面条件として PPE に入れる．
  \item \textbf{面演算子整合問題}：
        既知ジャンプを滑らかな補助圧力場へ埋め込まず，
        $G_\Gamma(p;j)=G(p)-B(j)$ のジャンプ補正付き面勾配で扱う．
  \item \textbf{高次ステンシル問題}：
        界面を跨ぐ CCD/FCCD ステンシルを禁止し，HFE で片側 Hermite データを延長する．
  \item \textbf{求解コスト問題}：
        高次演算子を直接「解く」ことに固執せず，欠陥補正で高次評価演算子と低次補正演算子を分離する．
\end{enumerate}
\end{tcolorbox}

\noindent\textbf{本章で定める閉包：}
結論として，本章が定める圧力ジャンプ型 PPE は
\[
  \text{相内定数密度 PPE}
  \;+\;
  G_\Gamma(p;j)=G(p)-B(j)
  \;+\;
  \mathrm{HFE}
  \;+\;
  \mathrm{DC}\ k=3
  \;+\;
  \text{大域ゲージ・再投影防護}
\]
である．FVM/FD，CSF，一括変密度 PPE，CCD-LU，
および滑らかな補助圧力場によるジャンプ表現は，
役割を限定して扱う数理モデルであり，本章の圧力ジャンプ型閉包そのものではない．
精度の全体像と閉包条件は §~\ref{sec:pressure_accuracy_summary} にまとめる．

%% ═══════════════════════════════════════════════════════════════════════════════
\subsection{CCD 法の楕円型求解法としての双対的役割}
\label{sec:ccd_elliptic}
%% ═══════════════════════════════════════════════════════════════════════════════

CCD 法は「既知の関数 $f$ の微分値 $f', f''$ を高精度に求めるスキーム」（役割1）に加えて，
\textbf{楕円型方程式の陰的求解法}（役割2）としても機能する．

\begin{itemize}
  \item \textbf{役割1（微分演算子）：} 既知の $f$（$\phi$，$\bu$ など）から $\Ord{h^6}$ 精度の
        $f', f'', f_{xy}$ を求める．曲率計算・粘性項評価に使用．
  \item \textbf{役割2（楕円型求解法）：} 未知の圧力 $p$ に対して CCD を適用し，
        $\Ord{h^6}$ 精度の演算子 $\mathbf{L}_{\mathrm{CCD}}^{\rho}$ を構成する（§~\ref{sec:ccd_poisson_matrix}）．
\end{itemize}

数学的構造の詳細（射影行列 $\mathbf{E}_1, \mathbf{E}_2$，演算子行列 $\mathbf{L}_{\mathrm{CCD}}^{\rho}$ の定義，
変密度 Poisson 演算子の導出）は付録~\ref{app:ccd_elliptic} を参照されたい．

%% ═══════════════════════════════════════════════════════════════════════════════
\subsection{CCD-Poisson 演算子の行列構造と変密度拡張}
\label{sec:ccd_poisson_matrix}
%% ═══════════════════════════════════════════════════════════════════════════════

§~\ref{sec:projection_derivation} で得た PPE（式~\eqref{eq:PPE_varrho}）を，
CCD 法（第~\ref{sec:CCD}章）を用いて $\Ord{h^6}$ 精度で離散化する．
本節では2つの課題を同時に解決する：

\begin{enumerate}
  \item \textbf{$\Ord{h^2} \to \Ord{h^6}$ への精度向上：}
        圧力 $p$，1次導関数 $p'$，2次導関数 $p''$ を連立未知量として同時求解する
        CCD ブロック系（付録~\ref{app:ccd_elliptic}）により，
        速度補正の圧力勾配を $\Ord{h^6}$ 精度で評価する
        （均一格子；非一様格子では面平均勾配 $\mathcal{G}^\text{adj}$ に切り替え，
        §~\ref{sec:fvm_ccd_corrector} 参照）．
  \item \textbf{チェッカーボード（奇偶デカップリング）の抑制：}
        コロケート配置では格子スケール振動が生じる．
        本稿では Balanced--Force 演算子整合（§~\ref{sec:balanced_force}）と
        FCCD 面ジェット部分系（§~\ref{sec:fccd_bf_sub_system}）の組み合わせにより
        $2\Delta x$ モードの発散駆動を抑制する（§~\ref{sec:checkerboard_solution}）．
\end{enumerate}

行列の代数的導出（$\mathbf{L}_{\mathrm{CCD}}^{(x)} = \mathbf{E}_2\,\mathcal{M}_{\mathrm{CCD}}^{-1}\,\mathcal{R}$，
演算子スペクトル半径 $\sigma_\mathrm{CCD}\,h^2 = 9.6$；§~\ref{sec:verify_ccd_spectral_radius}）は付録~\ref{app:ccd_elliptic}，
Neumann 境界条件の組み込み手順は付録~\ref{sec:ccd_neumann_bc} を参照のこと．
（$\sigma_\mathrm{CCD}$ は演算子スペクトル半径であり，密度 $\rho$ とは別記号．）

\subsubsection{2 次元変密度 CCD-Poisson 演算子の完全定式化}
\label{sec:ccd_2d_full}

2次元変密度 PPE $\bnabla\cdot(1/\rho\,\bnabla p) = q$ の CCD 離散化を示す．
$x$ 方向の演算子（式~\eqref{eq:L_split}）を格子点 $(i,j)$ で評価すると：
\begin{align}
  \bigl(\mathcal{L}_{\mathrm{CCD}}^{\rho,x}\bm{p}\bigr)_{i,j}
  &= \frac{1}{\rho_{i,j}}\bigl(D_x^{(2)} p\bigr)_{i,j}
  - \frac{(D_x^{(1)}\rho)_{i,j}}{\rho_{i,j}^2}\bigl(D_x^{(1)} p\bigr)_{i,j}
  \label{eq:Lx_varrho_discrete}
\end{align}
$y$ 方向も同様に：
\begin{align}
  \bigl(\mathcal{L}_{\mathrm{CCD}}^{\rho,y}\bm{p}\bigr)_{i,j}
  &= \frac{1}{\rho_{i,j}}\bigl(D_y^{(2)} p\bigr)_{i,j}
  - \frac{(D_y^{(1)}\rho)_{i,j}}{\rho_{i,j}^2}\bigl(D_y^{(1)} p\bigr)_{i,j}
  \label{eq:Ly_varrho_discrete}
\end{align}

合計演算子：
\begin{equation}
  \bigl(\mathcal{L}_{\mathrm{CCD}}^{\rho}\bm{p}\bigr)_{i,j}
  = \bigl(\mathcal{L}_{\mathrm{CCD}}^{\rho,x}\bm{p}\bigr)_{i,j}
  + \bigl(\mathcal{L}_{\mathrm{CCD}}^{\rho,y}\bm{p}\bigr)_{i,j}
  \label{eq:L_CCD_2d_full}
\end{equation}

\noindent\textbf{注：分相 PPE における第 2 項消失．}
式~\eqref{eq:Lx_varrho_discrete}--\eqref{eq:Ly_varrho_discrete} の
第 2 項（積の微分項）は $D^{(1)}\rho$ を含むため，
分相 PPE（§~\ref{sec:split_ppe_motivation}）では各相内で $\rho=\mathrm{const}$
となり\textbf{この項は厳密に零}である．
smoothed-Heaviside 一括 PPE でのみ第 2 項が界面帯で非零値を取る．

行列非形成評価手順（各反復で Thomas 法を $x$・$y$ 方向に交互適用，$\Ord{N_x N_y}$）は
付録~\ref{sec:ccd_matrix_free_eval} を参照のこと．

\subsubsection{Balanced--Force 原理から圧力ジャンプ型閉包への選別}
\label{sec:ccd_balanced_force}

§~\ref{sec:balanced_force} の本質は「CSF を使うこと」ではなく，
圧力勾配と界面応力を\textbf{同じ面・同じ係数・同じ発散演算子}で評価することである．
CSF 体積力
$\bm f_\sigma=\sigma\kappa_{lg}\bnabla\psi$ はこの原理を説明する
最小モデルとして有用だが，圧力ジャンプ型閉包としては二つの欠点を残す：
有限幅 $\delta_\varepsilon$ による $\Ord{\varepsilon^2}$ モデル誤差と，
圧力項とは別の体積力として入ることによる演算子位置のずれである．

したがって本章が採用するのは，CSF 体積力そのものではなく，
Balanced--Force の演算子位置一致要請を圧力ジャンプ形式へ移した
次の閉包である：
\[
  j_{gl}=p_g-p_l=-\sigma\kappa_{lg},
  \qquad
  G_{\Gamma,f}(p;j_{gl})=G_f(p)-B_f(j_{gl}).
\]
この形では，表面張力は予測段階右辺の体積力ではなく PPE と速度補正が共有する
面勾配契約として入る．
つまり，解くべき問題は「CSF を高次化すること」ではなく，
Young--Laplace ジャンプを投影演算子から消えない形で同じ面フラックスに載せることである．

\begin{tcolorbox}[resultbox, title={§9 の圧力閉包の選択}]
\begin{center}
\small
\renewcommand{\arraystretch}{\PaperTableArrayStretch}
\begin{tabular}{p{0.20\linewidth}p{0.36\linewidth}p{0.34\linewidth}}
\hline
候補 & 何を解こうとしたか & §9 の判定 \\
\hline
CSF + Balanced--Force & 圧力勾配と表面張力体積力の評価位置ずれを減らす & 原理の動機付け．最終的な毛管閉包ではない \\
滑らかな補助圧力場によるジャンプ表現 & 圧力ジャンプを連続場 $J$ として合成する & 楕円求解が $J$ を自由圧力として吸収し得るため採用しない \\
\textbf{ジャンプ補正付き面勾配} & 既知ジャンプを面勾配 $G_\Gamma=G-B$ に直接入れる & 採用．PPE と速度補正が同じ面条件を共有する \\
\hline
\end{tabular}
\end{center}

\smallskip
\noindent
HFE（§~\ref{sec:field_extension}）が必要になるのは，圧力ジャンプ型閉包が圧力場に
実ジャンプを持たせるからである．CSF のような滑らかな圧力場では HFE は本質ではないが，
圧力ジャンプ型分相 PPE では CCD/FCCD ステンシルが不連続値を跨がないための
必須部品になる．
\end{tcolorbox}

\subsubsection{空間離散化スキームの構造比較（FVM 基準離散化 vs CCD）}
\label{sec:ccd_vs_fvm}

表~\ref{tab:ccd_vs_fvm} は PPE の\textbf{空間離散化スキーム}（FVM 基準離散化 vs CCD）を比較する．
判断軸は空間精度と面演算子整合であり，FCCD/CCD を圧力閉包の
面演算子として採用する根拠を示す
（CCD 系を与件とした\textbf{線形求解法}の選択は §~\ref{sec:dc_comparison} で独立に扱う）．

\begin{table}[ht]
\centering
\caption{PPE 空間離散化スキームの構造比較}
\label{tab:ccd_vs_fvm}
\PaperCaptionNote{$N \times N$ 格子，$h = 1/N$ の設定で比較する．}
\PaperTableDenseSetup
\begin{tabularx}{\linewidth}{@{}p{0.19\linewidth}p{0.13\linewidth}X p{0.15\linewidth}@{}}
\toprule
空間離散化 & 空間精度 & 面演算子整合上の主な残差 & 本章での位置づけ \\
\midrule
FVM 基準離散化（2次中心差分） & $\Ord{h^2}$ & 低次フラックス閉包が律速 & 基準離散化 \\
\textbf{FCCD/CCD 面演算子}（本稿） & $\boldsymbol{\Ord{h^6}}$（平滑領域） & ジャンプ補正付き面勾配と同じ面で圧力補正を評価できる & 本章の面演算子 \\
\bottomrule
\end{tabularx}

\smallskip
注：寄生流れ量の数値評価は後続の静止液滴検証で扱う．
線形求解法の選択は §~\ref{sec:dc_comparison} を参照．
\end{table}
